\begin{table}[htbp]
\centering
\caption{Covariates included in the primary spatial model}
\label{tab:covariates}
\small\setlength{\tabcolsep}{3pt}
\revisioncolor
\begin{tabularx}{\linewidth}{@{}Xll@{}}
\toprule
Covariate & Analysis variable & Resolution\\
\midrule
Mean daily temperature ($^\circ$C) & \texttt{avg\_temp} & Monthly\\
Mean daily sunshine (h/day) & \texttt{sun\_time} & Monthly\\
Low-wind equivalent days & \texttt{stagnant\_days} & Monthly\\
Mean wind speed (m/s) & \texttt{avg\_wind} & Monthly\\
Mean relative humidity (\%) & \texttt{avg\_humid} & Monthly\\
Mean daily precipitation (mm/day) & \texttt{avg\_rain} & Monthly\\
Average daily distance per vehicle & \texttt{daily\_km} & Annual\\
Registered cars per capita & \texttt{cars\_per\_capita} & Annual\\
Population density & \texttt{pop\_density} & Annual\\
Industrial area per capita (m$^2$/person) & \texttt{industrial\_area} & Annual\\
Commercial area per capita (m$^2$/person) & \texttt{commercial\_area} & Annual\\
Green area per capita (m$^2$/person) & \texttt{green\_area\_per\_capita} & Annual\\
\bottomrule
\end{tabularx}
\begin{flushleft}\footnotesize\revisioncolor
Notes: Annual values repeat within district-year. All listed variables and their spatial lags enter Equation~(\ref{eq:sdm}). Weather labels describe station-level summaries before spatial interpolation; district low-wind-day values can be fractional. Coverage and processing qualifications are given in Supplementary Section S6.
\end{flushleft}
\end{table}
